\section{有機・高分子}

\subsection{平成25年度 問16}
\label{q:h25-reg-applied-16}

\noindent\textbf{原典：} sources/H25/17\_applied\_science.pdf、PDF物理7頁。

\textbf{公式正答：} 3（single）。

\par\textit{原典本文の校正と独自解説は未作成。分類のみの登録項目。}

\section{分析・分離}

\subsection{平成24年度 問16}
\label{q:h24-reg-applied-16}

\noindent\textbf{原典：} sources/H24/17\_applied\_science.pdf、PDF物理5頁。

\textbf{公式正答：} 5（single）。

\par\textit{原典本文の校正と独自解説は未作成。分類のみの登録項目。}

\subsection{平成24年度 問27}
\label{q:h24-reg-applied-27}

\noindent\textbf{原典：} sources/H24/17\_applied\_science.pdf、PDF物理8頁。

\textbf{公式正答：} 4（single）。

\par\textit{原典本文の校正と独自解説は未作成。分類のみの登録項目。}

\subsection{平成25年度 問15}
\label{q:h25-reg-applied-15}

\noindent\textbf{原典：} sources/H25/17\_applied\_science.pdf、PDF物理6頁。

\textbf{公式正答：} 1（single）。

\par\textit{原典本文の校正と独自解説は未作成。分類のみの登録項目。}

\subsection{平成28年度 問16}
\label{q:h28-reg-applied-16}

\noindent\textbf{原典：} sources/H28/17\_applied\_science.pdf、PDF物理6頁。

\textbf{公式正答：} 1（single）。

\par\textit{原典本文の校正と独自解説は未作成。分類のみの登録項目。}

\subsection{平成30年度 問13}
\label{q:h30-reg-applied-13}

\noindent\textbf{原典：} sources/H30/17\_applied\_science.pdf、PDF物理4頁。

\textbf{公式正答：} 2（single）。

\par\textit{原典本文の校正と独自解説は未作成。分類のみの登録項目。}

\subsection{平成30年度 問18}
\label{q:h30-reg-applied-18}

\noindent\textbf{原典：} sources/H30/17\_applied\_science.pdf、PDF物理6頁。

\textbf{公式正答：} 1（single）。

\par\textit{原典本文の校正と独自解説は未作成。分類のみの登録項目。}

\subsection{令和元年度 問12}
\label{q:r01-reg-applied-12}

\noindent\textbf{原典：} sources/R01/17\_applied\_science.pdf、PDF物理4頁。

\textbf{公式正答：} 4（single）。

\par\textit{原典本文の校正と独自解説は未作成。分類のみの登録項目。}

\subsection{令和2年度 問11}
\label{q:r02-reg-applied-11}

\noindent\textbf{原典：} sources/R02/17\_applied\_science.pdf、PDF物理4頁。

\textbf{公式正答：} 5（single）。

\par\textit{原典本文の校正と独自解説は未作成。分類のみの登録項目。}

\subsection{令和2年度 問17}
\label{q:r02-reg-applied-17}

\noindent\textbf{原典：} sources/R02/17\_applied\_science.pdf、PDF物理6頁。

\textbf{公式正答：} 4（single）。

\par\textit{原典本文の校正と独自解説は未作成。分類のみの登録項目。}

\subsection{令和3年度 問16}
\label{q:r03-reg-applied-16}

\noindent\textbf{原典：} sources/R03/17\_applied\_science.pdf、PDF物理6頁。

\textbf{公式正答：} 3（single）。

\par\textit{原典本文の校正と独自解説は未作成。分類のみの登録項目。}

\subsection{令和5年度 問15}
\label{q:r05-reg-applied-15}

\noindent\textbf{原典：} sources/R05/17\_applied\_science.pdf、PDF物理5頁。

\textbf{公式正答：} 2（single）。

\par\textit{原典本文の校正と独自解説は未作成。分類のみの登録項目。}

\subsection{令和6年度 問14}
\label{q:r06-reg-applied-14}

\noindent\textbf{原典：} sources/R06/17\_applied\_science.pdf、PDF物理4頁。

\textbf{公式正答：} 4（single）。

\par\textit{原典本文の校正と独自解説は未作成。分類のみの登録項目。}

\subsection{令和6年度 問19}
\label{q:r06-reg-applied-19}

\noindent\textbf{原典：} sources/R06/17\_applied\_science.pdf、PDF物理6頁。

\textbf{公式正答：} 1（single）。

\par\textit{原典本文の校正と独自解説は未作成。分類のみの登録項目。}

\subsection{令和7年度 問14}
\label{q:r07-reg-applied-14}

\noindent\textbf{原典：} sources/R07/17\_applied\_science.pdf、PDF物理5頁。

\textbf{公式正答：} 4（single）。

\par\textit{原典本文の校正と独自解説は未作成。分類のみの登録項目。}

\section{陽イオンの系統分析}

\subsection{平成26年度 問16}
\label{q:h26-reg-applied-16}

\noindent\textbf{原典：} sources/H26/17\_applied\_science.pdf、PDF物理7頁。

\textbf{公式正答：} 4（single）。

\paragraph{問題の要旨} 塩化物、酸性下の硫化物、水酸化物を順に除いたろ液から、最後に硫化物として分離する陽イオンを特定する。

\Needspace{4\baselineskip}\paragraph{確認した条件・記述（独自要約）}\mbox{}\par

\noindent\textbf{I：} 塩酸を滴下し沈殿をろ別\par

\noindent\textbf{II：} Iのろ液中でチオアセトアミドCH3CSNH2の加水分解を行い、生成した沈殿をろ別\par

\noindent\textbf{III：} IIのろ液中へアンモニア水を滴下し、生成した沈殿をろ別\par

\noindent\textbf{IV：} IIIのろ液中へ硫化水素を滴下し、生成した沈殿をろ別\par

\paragraph{問われる核心} 分析化学

\paragraph{必要知識} 系統的定性分析、溶解度積、硫化水素の酸解離、pHと沈殿選択性

\paragraph{考え方（査読待ち初稿）} 手順IでAg+をAgClとして除く。手順IIの酸性条件でCu2+が難溶性の硫化物となる。手順IIIでFe3+を水酸化物として除く。アンモニア性のろ液に硫化水素を加える手順IVではZn2+をZnSとして分離するので④。Na+はこれらの分離後も水溶液に残る。これは出題された古典的系統分析の条件に基づく説明で、実験条件一般への無条件な適用はしない。

\paragraph{他の選択肢（査読待ち初稿）} ①Na+はこの操作列での沈殿対象でない。②Ag+はI、③Cu2+はII、⑤Fe3+はIIIで先に除かれる。IとIVの取り違えは正答を変えるため、原典のローマ数字を必ず確認する。

\par 参考：\cite{qualitative_stanford}。

\par 難易度 3（校準初稿）、優先度未評価。

\section{未確定}

\subsection{平成29年度 問12}
\label{q:h29-reg-applied-12}

\noindent\textbf{原典：} sources/H29/17\_applied\_science.pdf、PDF物理4頁。

\textbf{公式正答：} 5（single）。

\par\textit{原典本文の校正と独自解説は未作成。分類のみの登録項目。}

\subsection{平成30年度 問12}
\label{q:h30-reg-applied-12}

\noindent\textbf{原典：} sources/H30/17\_applied\_science.pdf、PDF物理4頁。

\textbf{公式正答：} 5（single）。

\par\textit{原典本文の校正と独自解説は未作成。分類のみの登録項目。}

\subsection{令和2年度 問12}
\label{q:r02-reg-applied-12}

\noindent\textbf{原典：} sources/R02/17\_applied\_science.pdf、PDF物理4頁。

\textbf{公式正答：} 3（single）。

\par\textit{原典本文の校正と独自解説は未作成。分類のみの登録項目。}

\subsection{令和5年度 問22}
\label{q:r05-reg-applied-22}

\noindent\textbf{原典：} sources/R05/17\_applied\_science.pdf、PDF物理8頁。

\textbf{公式正答：} 4（single）。

\par\textit{原典本文の校正と独自解説は未作成。分類のみの登録項目。}
